% from ega1-1-fr.tex
\oldpage[I]{94}
En outre, ces homomorphismes satisfont aux conditions de compatibilité
usuelles~: pour $D(f')\supset D(g')$, le diagramme
\[
  \xymatrix{
    \Gamma(D(f'),\widetilde{A'})\ar[r]\ar[d] &
    \Gamma(D(f'),{}^a\varphi_*(\widetilde{A}))\ar[d]\\
    \Gamma(D(g'),\widetilde{A'})\ar[r] &
    \Gamma(D(g'),{}^a\varphi_*(\widetilde{A}))
  }
\]
est commutatif (\hyperref[0.1.5.1-fr]{0, 1.5.1})~; on a donc bien défini
un homomorphisme de $\mathscr{O}_{X'}$-Algèbres, les $D(f')$ formant une
base de la topologie de $X'$ (\hyperref[0.3.2.3-fr]{0, 3.2.3}). Le couple
$\Phi=({}^a\varphi,\widetilde{\varphi})$ est donc un \emph{morphisme}
d'espaces annelés
\[
  \Phi:(X,\mathscr{O}_X)\to(X',\mathscr{O}_{X'})
\]
(\hyperref[0.4.1.1-fr]{0, 4.1.1}).

Notons en outre que, si on pose $x'={}^a\varphi(x)$, l'homomorphisme
$\widetilde{\varphi}_x^\#$ (\hyperref[0.3.7.1-fr]{0, 3.7.1}) n'est autre
que l'homomorphisme
\[
  \varphi_x:A'_{x'}\to A_x
\]
déduit canoniquement de $\varphi:A'\to A$
(\hyperref[0.1.5.1-fr]{0, 1.5.1}). En effet, tout $z'\in A'_{x'}$
s'écrit $g'/f'$, où $f'$, $g'$ sont dans $A'$ et
$f'\notin\mathfrak{j}_{x'}$~; $D(f')$ est donc un voisinage de $x'$ dans
$X'$, et l'homomorphisme
$\Gamma(D(f'),\widetilde{A'})\to
\Gamma({}^a\varphi^{-1}(D(f')),\widetilde{A})$ déduit de
$\widetilde{\varphi}$ n'est autre que $\varphi_{f'}$~; en considérant la
section $s'\in\Gamma(D(f'),\widetilde{A'})$ correspondant à
$g'/f'\in A'_{f'}$, on obtient bien
$\widetilde{\varphi}_x^\#(z')=\varphi(g')/\varphi(f')$ dans $A_x$.
\end{env}

\begin{example}[1.6.2]
\label{I.1.6.2-fr}
Soient $S$ une partie multiplicative de $A$, $\varphi$ l'homomorphisme
canonique $A\to S^{-1}A$~; on a déjà vu
(\hyperref[I.1.2.6-fr]{1.2.6}) que ${}^a\varphi$ est un
\emph{homéomorphisme} de $Y=\operatorname{Spec}(S^{-1}A)$ sur le
sous-espace de $X=\operatorname{Spec}(A)$ formé des $x$ tels que
$\mathfrak{j}_x\cap S=\varnothing$. En outre, pour tout $x$ appartenant à
ce sous-espace, donc de la forme ${}^a\varphi(y)$ avec $y\in Y$,
l'homomorphisme
$\widetilde{\varphi}_y^\#:\mathscr{O}_x\to\mathscr{O}_y$ est
\emph{bijectif} (\hyperref[0.1.2.6-fr]{0, 1.2.6})~; autrement dit,
$\mathscr{O}_Y$ s'identifie au faisceau induit sur $Y$ par $\mathscr{O}_X$.
\end{example}

\begin{proposition}[1.6.3]
\label{I.1.6.3-fr}
Pour tout $A$-module $M$, il existe un isomorphisme canonique fonctoriel du
$\mathscr{O}_{X'}$-Module $(M_{[\varphi]})^{\sim}$ sur l'image directe
$\Phi_*(\widetilde{M})$.
\end{proposition}

Posons pour abréger $M'=M_{[\varphi]}$, et pour tout $f'\in A'$, posons
encore $f=\varphi(f')$. Les modules de sections
$\Gamma(D(f'),\widetilde{M'})$ et $\Gamma(D(f),\widetilde{M})$
s'identifient respectivement aux modules $M'_{f'}$ et $M_f$ (sur
$A'_{f'}$ et $A_f$ respectivement)~; en outre, le $A'_{f'}$-module
$(M_f)_{[\varphi_{f'}]}$ est canoniquement isomorphe à $M'_{f'}$
(\hyperref[0.1.5.2-fr]{0, 1.5.2}). On a donc un isomorphisme fonctoriel de
$\Gamma(D(f'),\widetilde{A'})$-Modules~:
$\Gamma(D(f'),\widetilde{M'})\xrightarrow{\sim}
\Gamma({}^a\varphi^{-1}(D(f')),\widetilde{M})_{[\varphi_{f'}]}$ et ces
isomorphismes satisfont aux conditions de compatibilité usuelles avec les
restrictions (\hyperref[0.1.5.6-fr]{0, 1.5.6}), donc définissent
l'isomorphisme fonctoriel annoncé. On notera que, de façon précise, si
$u:M_1\to M_2$ est un homomorphisme de $A$-modules, il peut être considéré
comme un homomorphisme $(M_1)_{[\varphi]}\to(M_2)_{[\varphi]}$ de
$A'$-modules~; si on note $u_{[\varphi]}$ cet homomorphisme,
$\Phi_*(\widetilde{u})$ s'identifie à $(u_{[\varphi]})^{\sim}$.

Cette démonstration prouve aussi que pour toute $A$-algèbre $B$,
l'isomorphisme
